\section{Architecture of the formalization}
\label{sec:architecture}

\subsection{Scale and shape}

Table~\ref{tab:scale} gives the size of the repository. The two developments are
deliberately kept in separate libraries with separate roots, so that the Euler material
cannot depend on the Navier--Stokes material or conversely.

\begin{table}[t]
\centering
\small
\begin{tabular}{@{}lrrr@{}}
\toprule
library & \Lean{} files & lines & \lean{theorem}/\lean{lemma} \\
\midrule
\file{NavierStokes} & 753 & 405{,}375 & 24{,}694 \\
\file{Euler} & 1{,}829 & 212{,}155 & 10{,}594 \\
\file{ComparatorChallenges} & 2 & 474 & 9 \\
\midrule
total & 2{,}584 & 618{,}004 & 35{,}297 \\
\bottomrule
\end{tabular}
\caption{Size of the development after the dead-code removal of
Section~\ref{sec:cleanup}, counted over the library directories
(\file{NavierStokes.lean} and \file{Euler.lean} add two root files, for 2{,}588 in the tree). Lines are physical
lines of Lean source including comments.}
\label{tab:scale}
\end{table}

The import graph is unusually \emph{flat}. On average a file imports 2.74 modules, of
which 2.30 are project-local; 73 files import no project-local module at all (they sit
directly on \Mathlib); the longest chain of project-local imports is 96 modules; and only
57 modules are imported by eight or more others. The largest single files are
\file{Euler/EulerProof.lean} (20{,}755 lines), \file{NavierStokes/CorrectionStep.lean}
(9{,}849), \file{NavierStokes/CorrectionInitializationNoOptions.lean} (5{,}573) and
\file{NavierStokes/CorrectionInitialization.lean} (5{,}562); everything else is under
3{,}000 lines, ninety per cent of them are under 600 lines, and the median file is 114
lines long.

This shape is a consequence of how the development was produced and is worth commenting
on for anyone planning a large formalization. Each file introduces a small number of
definitions and proves the lemmas about them; because the files are small, elaboration is
parallel, failures are local, and a reader can find the place where a specific estimate is
proved by searching for the constant it mentions. The price is that there is almost no
reuse: the same argument is often repeated in two files that cannot see each other,
because the module that would have shared it would have to be an ancestor of both. We
measured the extreme case: 74 of the 2,659 modules are not reachable from any of the
repository's entry points (Section~\ref{sec:cleanup}).

\subsection{Three layers}
\label{sec:layers}

It is useful to separate three layers, because they are checked against different
standards.

\paragraph{The statement layer.} A handful of modules contain no analysis at all:
\file{NavierStokes/ProblemStatement.lean} (the torus statement),
\file{NavierStokes/R3/ProblemStatement.lean} (the whole-space statement),
\file{Euler/SolutionDefinitions.lean} and \file{ComparatorChallenges/*.lean}. These are
where fidelity to the source papers and to the Millennium problem statement is decided.
They define the carriers, the differential operators \eqref{eq:residual} and the
structures \lean{CandidateProperties}, \lean{GlobalFiniteEnergySolution},
\lean{GlobalSmoothSolution} and \lean{InitialVelocityConditionDecay} that spell out what
``a solution'' is. Everything in this layer is short ($\sim$150 lines each) and is meant to
be read by a referee.

\paragraph{The infrastructure layer.} About 3{,}800 lines in 18 modules supply analysis
that \Mathlib{} does not have: smooth cutoffs with all-order derivative bounds, a
one-sided exponential edge $\exp(-c/x)$ with its inverse-power family, zero extension
across such an edge, finite-order jet bookkeeping in weighted band-indexed classes, the
Banach space $C^k([-1,1];E)$ as a closed submodule of jet data, all-orders differentiation
under the integral sign, nested $L^2$ as a product $L^2$, mollification on $L^2$, Sobolev
spaces on a cylinder as a word lattice, and a heat-kernel gain of derivatives. These
modules are the natural candidates for replacement by a general library, and they are the
subject of Section~\ref{sec:libraries}.

\paragraph{The construction layer.} Everything else: the selected profiles, the geometric
regions, the dyadic charts, the wave labels, the correction stages, the packet estimates,
and the assembly theorems that instantiate the statement layer. This is where the
mathematics lives, and it is where a reader will spend their time.

\subsection{The dependency chain of a headline theorem}
\label{sec:chain}

The two Navier--Stokes alternatives are proved along the following path. The whole-space
statement is the primitive one; the periodic corollary is derived from it.

\begin{center}
\begin{tikzpicture}[node distance=5mm and 7mm,
   box/.style={draw,rounded corners=1pt,fill=blue!6,align=center,font=\scriptsize,inner sep=3pt,text width=3.05cm},
   arr/.style={-{Latex[length=2mm]},draw=black!70}]
 \node[box] (ps) {\file{ProblemStatement}\\operators, residual};
 \node[box,right=of ps] (r3ps) {\file{R3/ProblemStatement}\\whole-space structures};
 \node[box,right=of r3ps] (act) {\file{R3/ActualCandidate}\\selected construction};
 \node[box,below=of act] (brk) {\file{R3/CandidateBreakdown}\\no global competitor};
 \node[box,left=of brk] (vis) {\file{R3/ViscosityScaling}\\$\nu$-scaling};
 \node[box,left=of vis] (thm) {\file{R3/Theorem}\\Theorem~\ref{thm:ns-r3}};
 \node[box,below=of thm] (par) {\file{R3/ParabolicScaling}\\compression};
 \node[box,right=of par] (pp) {\file{PeriodizePDE}\\periodic lift};
 \node[box,right=of pp] (pc) {\file{PeriodicPaperTheorem}\\Theorem~\ref{thm:ns-torus}};
 \node[box,below=of thm,yshift=-11mm] (cmp) {\file{ComparatorSolution}\\(C) and (D)};
 \draw[arr] (ps)--(r3ps); \draw[arr] (r3ps)--(act); \draw[arr] (act)--(brk);
 \draw[arr] (brk)--(vis); \draw[arr] (vis)--(thm);
 \draw[arr] (thm)--(par); \draw[arr] (par)--(pp); \draw[arr] (pp)--(pc);
 \draw[arr] (pc.south) -- ++(0,-3mm) -| (cmp.east);
 \draw[arr] (thm.south) -- ++(0,-3mm) -| ([xshift=6mm]cmp.north);
\end{tikzpicture}
\end{center}

The Euler chain is different in kind: the non-continuation statement
(Theorem~\ref{thm:euler-bkm}) is a corollary of the quantitative one
(Theorem~\ref{thm:euler-quant}) through the Beale--Kato--Majda criterion, and the
quantitative statement is assembled from the iteration.

\subsection{Naming and documentation}
\label{sec:naming}

Two conventions carried the development and are worth recording.

\emph{Names say what is being proved, not where it is used.}
There are modules called \file{ReservedPatches}, \file{ExponentLedger},
\file{FiveProfileMoments}, \file{MeanLocalDefectBounds}, \file{OutgoingEntranceCone},
\file{SignedRequestContinuation}. A reader who knows the source paper can find the Lean
counterpart of a numbered proposition by looking for the object it is about. This is the
opposite of a ``Chapter3\_Lemma7'' convention, and it is the right choice at this scale.

\emph{Module docstrings state the role, and often the caveat.}
\file{NavierStokes/ExponentLedger.lean} opens with ``The candidate manuscript,
Proposition 10.3, assigns real exponents to analytic estimates. This file checks the
arithmetic of those assignments, conditional on the estimates being valid. It does not
define the analytic classes, construct a correction, or prove an estimate for a PDE.''
A formalization whose docstrings say what they do \emph{not} prove is much easier to audit
than one whose docstrings only advertise.

\subsection{Dead code, removed}
\label{sec:cleanup}

Reachability analysis from the repository's entry points
(\file{NavierStokes.lean}, \file{Euler.lean}, \file{Euler/Solution.lean},
\file{ComparatorChallenges/NavierStokes.lean}, \file{ComparatorChallenges/Euler.lean},
\file{NavierStokes/TauCetiBridge.lean}) left 75 modules that no entry point reaches. 74 of
them were deleted during the migration---26{,}089 lines, 4.0\% of the tree:

\begin{center}\small
\begin{tabular}{@{}lrr@{}}
\toprule
 & files & lines \\
\midrule
before & 2{,}659 & 643{,}986 \\
removed (\lean{*NoOptions*} copies of existing modules, in a separate namespace) & 11 & 7{,}032 \\
removed (investigations and orphaned fragments of earlier constructions) & 63 & 19{,}057 \\
after & 2{,}585 & 617{,}967 \\
\bottomrule
\end{tabular}
\end{center}

The deletion is safe by construction rather than by inspection: no remaining file imports any
of the removed modules, which was checked after the removal (zero dangling imports), so no
reachable declaration changed meaning. It is also reversible with a single \lean{git revert}.
The one unreachable module that was kept is \file{NavierStokes/TauCetiBridge.lean}, which is
deliberately a leaf.

This is the only substantial reduction of code available in the repository, and it is worth
being precise about why. Removing dead code took out 4\% of the lines at zero risk. Replacing
project code by library code would take out 162 lines (the two ``free'' files of
Table~\ref{tab:infra}) plus part of nine others, and carries the risk of touching
declarations the unfinished build depends on. The remaining 99\% of the code is not
duplicated library material and cannot be reduced by substitution; it is the proof.
